% ============================================================
% BÖLÜM 12: ÇOKLU AJAN DEBUG PIPELINE
% ============================================================
\chapter{Çoklu Ajan Debug Pipeline}
\label{ch:multi_agent}

\section{AI-Orchestrated R\&D Yaklaşımı}
\label{sec:ma_orchestrated}

Bu bölüm, INNATE projesinin yaklaşık üç ay süren geliştirme sürecinde
benimsenen çalışma tarzını sistematik olarak belgeler. Söz konusu tarz,
geleneksel ``tek mühendis tek editör başında kod yazar'' modelinden
belirgin biçimde ayrılır. Berke Tezgöçen'in projedeki rolü, klasik
anlamda bir kod yazarı olmaktan çok bir \emph{orkestra şefi}'ne benzer:
hangi enstrümanın (ajanın) ne zaman, hangi nüansla çalacağına karar
verir; doğrudan tuşlara basan kendisi değildir. Bu yaklaşımı kısaca
\textbf{AI-orchestrated mühendislik} olarak adlandıracağız.

\subsection{Berke'nin Çalışma Profili}

Berke, projenin başında kendisi için açık bir öz-değerlendirme yapmıştır:
\emph{``Ben kod yazmıyorum, sistem düşünüyorum.''} Bu cümle bir zayıflık
itirafı olarak değil, bir mühendislik pozisyonunun açık beyanı olarak
okunmalıdır. Pratikte bu, şu davranış kalıpları anlamına gelir:

\begin{itemize}
  \item \textbf{Üst seviye karar verme:} Hangi mimari (PINN, Fourier
        Neural Operator, fiziksel-modüler INNATE) hangi probleme uyar,
        hangi Reynolds aralığı eğitimde kullanılır, hangi veri seti
        bırakılır, hangi metrik birincil hedeftir gibi sorulara
        Berke karar verir.
  \item \textbf{Doğrulama ve sorgulama:} Bir ajanın ürettiği çözümü
        çoğu zaman doğrudan kabul etmez; sayısal sonuçları, terminal
        çıktısını, log dosyalarını okur, gerekirse ekran fotoğrafı ile
        ikinci bir ajana gösterir.
  \item \textbf{Ajan delegasyonu:} Kod yazımı, formül türetme, hata
        ayıklama, dokümantasyon gibi alt görevler farklı ajanlara
        dağıtılır. Hangi ajanın hangi görev için uygun olduğunu Berke
        seçer.
  \item \textbf{Sistem-seviyesi entegrasyon:} Birden fazla ajanın çelişen
        önerileri olduğunda nihai birleştirici karar yine Berke'dedir.
\end{itemize}

Bu profil, çağdaş yazılım/araştırma mühendisliğinde yükselen bir trende
karşılık gelir: yapay zekâ ajanlarının yalnızca bir ``yazım yardımcısı''
değil, birer \emph{çoğullaştırılmış araştırmacı} olarak kullanılması.
Bir kişi, eğer doğru sorular sorabiliyor ve doğru biçimde delege
edebiliyorsa, klasik anlamda tek başına asla başaramayacağı bir
çok-uzmanlıklı projeyi yürütebilir.

\subsection{Geleneksel Solo Geliştirme ile Karşılaştırma}

Tablo~\ref{tab:ma_solo_vs_multi}, klasik solo geliştirme ile bu projede
benimsenen çoklu ajan yaklaşımı arasındaki temel farkları özetler.

\begin{table}[H]
\centering
\small
\begin{tabular}{p{4cm} p{5cm} p{5cm}}
\toprule
\textbf{Boyut} & \textbf{Solo Geliştirme} & \textbf{Çoklu Ajan Pipeline} \\
\midrule
Bilgi kapsamı &
Geliştiricinin kişisel uzmanlığı ile sınırlı &
Türbülans, ML, sayısal analiz, sistem mühendisliği gibi farklı
alanlara paralel erişim \\
\addlinespace
Hata ayıklama hızı &
Hipotez $\rightarrow$ tek deneme $\rightarrow$ döngü &
Hipotez $\rightarrow$ paralel ikinci/üçüncü görüş $\rightarrow$
muhalif test $\rightarrow$ deneme \\
\addlinespace
Kör nokta riski &
Yüksek; kişisel önyargılar zor fark edilir &
Görece düşük; bağımsız ajanlar farklı önyargılar taşır \\
\addlinespace
Halüsinasyon riski &
Düşük (insan kasıtlı yanılmaz) ama insan yorgunluğu yüksek &
Yapay; ajan halüsine edebilir, doğrulama \emph{şart} \\
\addlinespace
Karar yetkisi &
Kişide &
Yine kişide; ajanlar danışmandır, karar mercii değil \\
\addlinespace
Belge bırakma alışkanlığı &
Çoğu zaman zayıf (zaman bulunamıyor) &
Yüksek; her ajan sohbeti zaten metinsel bir kayıttır \\
\bottomrule
\end{tabular}
\caption{Solo geliştirme ile çoklu ajan pipeline karşılaştırması.}
\label{tab:ma_solo_vs_multi}
\end{table}

\subsection{Yöntemin Kısıtları}

Bu yaklaşımı bir başarı reçetesi olarak sunmak yanıltıcı olur. Üç ay
boyunca açıkça gözlemlenen kısıtlar şunlardır:

\begin{enumerate}
  \item \textbf{Halüsinasyon:} Ajanlar, var olmayan dosyalardan,
        yanlış API isimlerinden veya yanlış matematiksel
        özdeşliklerden bahsedebilir. Bu nedenle her teknik iddia, en
        az bir bağımsız doğrulamadan geçirilmelidir.
  \item \textbf{Konsensüs paralizi:} Çok sayıda ajan aynı soruya
        cevap verdiğinde sonuçlar çoğu zaman birbirine yakın çıkar;
        bu yakınlık ``doğruluk'' ile karıştırılmamalıdır. Aynı veri
        seti üzerinde eğitilmiş modeller benzer hataları paylaşabilir.
  \item \textbf{Maliyet ve gecikme:} Yüksek akıl yürütme seviyesinde
        çağrılar (örn. xhigh) tek başına 7--10 dakika sürebilir;
        ajanın ajan içinden çağrılması bu süreyi 20+ dakikaya çıkarır.
        Bu mühendislik gerçekleri sürecin tasarımını biçimlendirir.
  \item \textbf{Hesap verebilirlik:} Bir hata yapıldığında ``ajan
        söyledi'' bir mazeret değildir. Nihai sorumluluk, kararı
        veren insandadır.
\end{enumerate}

\noindent
Bu kısıtlar, sonraki bölümlerde tarif edilen pipeline'ın \emph{neden}
bu biçimde tasarlandığını anlamak için temel arka plandır.

% ============================================================
\section{Kullanılan Ajan Takımı}
\label{sec:ma_takim}

INNATE projesinde aktif olarak görev almış ajanlar ve rolleri,
Tablo~\ref{tab:ma_ajanlar}'da özetlenmiştir. Tablo, her ajanın
yalnızca model adını değil, projedeki somut katkı örneğini de
içerir. Bu, soyut bir ``AI ekibi'' tanımı yerine, gerçek mühendislik
görevlerine bağlanmış bir takım resmi sunar.

\begin{table}[H]
\centering
\footnotesize
\begin{tabular}{p{2.6cm} p{2.6cm} p{3.4cm} p{4.4cm}}
\toprule
\textbf{Ajan} & \textbf{Model} & \textbf{Rol} & \textbf{Örnek katkı} \\
\midrule
Claude (ana koordinatör) &
Anthropic Opus 4.7 (1M ctx) &
Sistem entegrasyonu, kod yazma, akıl yürütme &
INNATE v2/v3 model dosyalarının yazımı, training pipeline kurulumu \\
\addlinespace
Codex CLI &
OpenAI GPT-5.5 (xhigh) &
Matematiksel türetme, ikinci görüş, algoritma doğrulama &
RK4-IMEX şemasının kararlılık çözümlemesi, von-Neumann sınırları \\
\addlinespace
Gemini Pro &
Google Gemini 3.1 Pro &
Üçüncü görüş, kod review &
Spectral loss tasarımının bağımsız değerlendirmesi \\
\addlinespace
Cursor IDE Agent &
Hibrit (uygulama içi) &
Bağımsız bug-fix, IDE düzeyinde refaktör &
Phase~C divergence için $\nu_t$ cap önerisi (von-Neumann) \\
\addlinespace
CFD Expert &
Opus alt-ajanı &
Türbülans, LES, kapama modelleri &
Buoyancy damping'in katman başına bölünmesi gerektiğinin tespiti \\
\addlinespace
ML Expert &
Opus alt-ajanı &
PyTorch, optimizer, neural operator &
20-katmanlı fractional-step LES'in safety çarpımı uyarısı \\
\addlinespace
Reviewer &
Opus alt-ajanı &
Kod review, kalite kontrol &
INNATE kütüphanesi 6-bug fix'ten sonra patch doğrulaması \\
\addlinespace
Tenth Man &
Codex GPT-5.5 (xhigh) &
Zorunlu muhalif (devil's advocate) &
9 ajan ``MLP fc2 split'' fix'ini onayladığında karşı argümanı \\
\bottomrule
\end{tabular}
\caption{INNATE projesinde aktif olarak kullanılan ajan takımı.}
\label{tab:ma_ajanlar}
\end{table}

\subsection{Rollerin Mantığı}

Takım, rastgele bir AI yığını değildir; her ajan belirli bir epistemik
işlev için seçilmiştir.

\begin{itemize}
  \item \textbf{Ana koordinatör (Claude):} Uzun bağlam penceresi
        sayesinde proje tarihçesini, log dosyalarını ve kod tabanını
        bir arada tutar. Pratikte ``proje belleği'' işlevi görür.
  \item \textbf{Matematik motoru (Codex):} Türetme yoğun işlerde,
        özellikle kararlılık analizi, türev çıkarımı ve özel fonksiyon
        sınırlarında öne çıkar. Ana koordinatörün matematik kısa
        devrelerine karşı bir ``ikinci göz''dür.
  \item \textbf{Bağımsız üçüncü göz (Gemini):} Anthropic dışı bir
        sistem olarak Claude/Codex pipeline'ında ortak olabilecek
        eğitim verisi önyargılarını taşımaz. Üçüncü görüş kanalı bu
        nedenle değerlidir.
  \item \textbf{IDE-içi ajan (Cursor):} Ana pipeline'a bağımlı değildir;
        Berke'nin paralel olarak çalıştırdığı bir ortamdır. Bu
        bağımsızlık, çapraz-doğrulama açısından kıymetlidir
        (Bölüm~\ref{sec:ma_cursor}).
  \item \textbf{Uzmanlık alt-ajanları (CFD, ML, Reviewer):} Ana
        koordinatör tarafından çağrılan, dar uzmanlığa odaklı görev
        kanallarıdır. Pratikte ``konunun adamına sor'' davranışını
        otomatikleştirirler.
  \item \textbf{Onuncu Adam (Tenth Man):} Adını, dokuz danışmanın
        hemfikir olduğunda zorunlu olarak muhalif kalmakla görevli
        onuncu kişi prensibinden alır. Konsensüsün doğruluk ile
        karıştırılmasını önler.
\end{itemize}

% ============================================================
\section{Multi-Agent Review Pipeline}
\label{sec:ma_pipeline}

Projedeki tipik bir hata-ayıklama veya tasarım kararı, neredeyse her
zaman benzer bir dizilim üzerinden ilerlemiştir. Bu dizilim, zaman
içinde rafine olarak Şekil~\ref{fig:ma_pipeline_text}'teki kavramsal
akışa oturmuştur.

\begin{figure}[H]
\centering
\begin{tabular}{l}
\textbf{(1)} Berke gözlem yapar (terminal, log, ekran fotoğrafı). \\
\quad $\downarrow$ \\
\textbf{(2)} Claude bir tanı/hipotez üretir. \\
\quad $\downarrow$ \\
\textbf{(3)} Codex'e ikinci görüş: matematik doğrulama, alternatif hipotez. \\
\quad $\downarrow$ \\
\textbf{(4)} Gerekirse Gemini'ye üçüncü görüş gönderilir. \\
\quad $\downarrow$ \\
\textbf{(5)} Konsensüs varsa Tenth Man muhalif test'e tabi tutulur. \\
\quad $\downarrow$ \\
\textbf{(6)} Berke nihai kararı verir; uygulama Claude tarafından yapılır. \\
\quad $\downarrow$ \\
\textbf{(7)} Reviewer ajan, uygulanan değişikliği bağımsız okur ve raporlar. \\
\end{tabular}
\caption{Tipik çoklu ajan review pipeline'ı.}
\label{fig:ma_pipeline_text}
\end{figure}

Bu akış katı bir protokol değil, bir \emph{varsayılan davranış}tır.
Basit, düşük etkili değişikliklerde adım~3--5 atlanabilir. Mimari
düzeyde değişikliklerde ise pipeline'ın sonuna kadar gidilmesi
beklenir.

\subsection{Vaka Analizi 1 — Yedi BLOCKER Bug Toplu Fix'i (2026-05-01)}

01 Mayıs 2026'da, INNATE training pipeline'ının yeni bir sürümü için
çoklu ajan code review oturumu düzenlenmiştir. Yöntem şu şekildedir:

\begin{enumerate}
  \item Aynı kod tabanı, aynı talimatla \emph{paralel} olarak dört
        ayrı ajana review için verilmiştir: Codex, Gemini, ML Expert
        ve Reviewer.
  \item Her ajandan üç kategoride bulgu çıkarması istenmiştir:
        BLOCKER (yayına engel), MAJOR, MINOR.
  \item Aynı bulgu en az iki bağımsız ajan tarafından raporlanmışsa
        BLOCKER olarak kabul edilmiştir.
  \item Bu eşikleme sonucunda yedi BLOCKER bulgu konsolide edilmiştir.
\end{enumerate}

\noindent
Bu yöntemin değeri, ``hangi ajanın haklı olduğunu'' tartışmaktan
kaçınmasında yatar. Bağımsız ajanların aynı hatayı işaret etmesi,
hatanın gerçek olma olasılığını çarpıcı biçimde artırır. Ardından
yedi BLOCKER, tek bir tutarlı yamada toplu olarak düzeltilmiş, sonuç
yeniden review'a verilmiştir.

\subsection{Vaka Analizi 2 — Phase~C Divergence (2026-05-08)}

Eğitimin Phase~C aşamasında ısıl alan, $\theta_{\mathrm{rms}}$
metriği üzerinden patlama davranışı sergilemiştir. Akış sırası:

\begin{enumerate}
  \item \textbf{Gözlem:} $\theta_{\mathrm{rms}}$ Phase~B sonunda
        $\sim 0.30$~LES, Phase~C başında $\sim 1.17$~LES bandını
        aşarak patlama eğilimi gösterdi.
  \item \textbf{Claude'un tanısı:} Forcing UNFREEZE artığı bug;
        scale-similarity ve $C_s$ MLP'si ısıl döngüye geri besleniyor.
  \item \textbf{Codex konsültasyonu:} ``Stratejinin yönü doğru,
        derecesi yetersiz.'' Bu cümle, tek bir cümlede tüm vakanın
        karakterini özetler.
  \item \textbf{Önerilen dört dokunuş:}
        \begin{enumerate}
          \item Phase~C başlangıcında ``jump'' parametresinin sıfırlanması.
          \item Optimizer momentum buffer'larının resetlenmesi.
          \item Spektral kayıp tanımında $\log(\cdot + \varepsilon)$
                gibi sayısal güvenlik.
          \item $C_s$ MLP çıktısının fiziksel olarak makul bir aralığa
                kelepçelenmesi (bounding).
        \end{enumerate}
  \item \textbf{Karar:} Berke dört dokunuşu da uygulamayı seçti.
        Hiçbiri tek başına çözümün tamamı değildi; birleşimleri
        gradyanı yumuşatıp Phase~C'yi stabilize etti.
\end{enumerate}

Bu vakanın metodolojik dersi şudur: tek bir ``akıllı'' ajanın tek
``parlak'' fikrini aramak yerine, birden çok küçük dokunuşun
birleşimini test etmek çoğu zaman daha sağlam bir mühendislik
çözümüdür. Çoklu ajan pipeline'ı tam olarak bu birleşimi üretmeyi
kolaylaştırır.

% ============================================================
\section{Codex Konsültasyon Felsefesi}
\label{sec:ma_codex}

Pipeline'ın bel kemiği niteliğindeki bir bileşen, Codex'e gönderilen
``ikinci görüş'' konsültasyonlarıdır. Bu bölüm, neden böyle bir
disiplinin benimsendiğini açıklar.

\subsection{Neden Tek Ajan Yetmez?}

Tek bir dil modeline güvenmenin üç sistematik zayıflığı vardır:

\begin{itemize}
  \item \textbf{Halüsinasyon:} Model, içsel olarak ``emin'' hissetse
        bile, gerçekte var olmayan bir API, formül veya teorem üretebilir.
  \item \textbf{Dar görüş alanı:} Bağlam penceresi içine alınan
        bilgilere göre cevap verir; o pencereye girmeyen kritik bir
        ayrıntıyı kendiliğinden hatırlamak zorunda değildir.
  \item \textbf{Aşırı güven (over-confidence):} Akıcı dil, doğruluğun
        kendisi değildir; aksine, akıcı yanlış cevaplar daha tehlikelidir
        çünkü insanı ikna eder.
\end{itemize}

İkinci bir bağımsız sistem (Codex), bu üç zayıflığı birden zayıflatır.
İkinci ajanın aynı yanlışı paylaşma ihtimali, eğitildiği veri ve
mimari farklılıkları nedeniyle düşüktür.

\subsection{Çağrı Modları ve Maliyetleri}

Codex çağrıları, akıl yürütme efor seviyesine göre kademelenir:
\texttt{low}, \texttt{medium}, \texttt{high} ve \texttt{xhigh}.
İnsancıl bir analoji kurmak gerekirse:

\begin{itemize}
  \item \texttt{medium}: Bir doktora öğrencisinin günlük içgüdüsel cevabı.
  \item \texttt{high}: Aynı öğrencinin tahtaya çıkıp türettiği cevabı.
  \item \texttt{xhigh}: Aynı öğrencinin bir gece düşündükten sonra
        yazdığı cevabı.
\end{itemize}

\noindent
\texttt{xhigh} seviyesi, çağrı başına yaklaşık 7--10 dakika sürebilir;
buna karşılık matematiksel olarak en az hata yapan moddur. Pratik
gözlem şöyledir: \texttt{xhigh}, derin algoritma ve kararlılık
analizinde kullanılmalı; basit LaTeX/prompt üretimi için
\texttt{medium} fazlasıyla yeter. Bu ekonomi-akıl yürütme dengesi,
pipeline'ın ne kadar verimli çalışacağını doğrudan belirler.

\subsection{Ajan-içi Codex Çağrısı Tuzağı}

20 Nisan 2026'da yapılan bir deneyde, Codex \texttt{xhigh} bir alt
ajanın içinden çağrıldığında çağrı 20 dakikadan fazla sürdü ve fiili
olarak tamamlanmadı. Aynı görev, doğrudan Bash üzerinden Codex'e
verildiğinde 3--4 dakikada bitti. Bu fiyaskonun dersleri pipeline'ın
çalışma kurallarına yazıldı:

\begin{itemize}
  \item Doğrudan kullanım: \texttt{Bash(codex exec ...)}.
  \item Uzun teknik danışmanlık: özel olarak bu iş için tasarlanmış
        \texttt{codex-consultant} alt-ajanı.
  \item Yasaklanan yol: gelişigüzel bir başka ajanın içinden
        \texttt{codex exec} tetiklemek.
\end{itemize}

Bu kuralın anlamı, ``araç-içinde-araç'' soyutlamasının kâğıt üzerinde
zarif görünse de pratikte ek gecikme ve hata kaynağı yarattığıdır.

% ============================================================
\section{Cursor Ajanının Bağımsız Katkısı}
\label{sec:ma_cursor}

Berke, pipeline'ı yalnızca Claude/Codex ekseninde işletmemiştir.
Paralel olarak Cursor IDE içinde de bir ajan çalıştırmıştır. Bu
ikinci ortam, ana pipeline'a bağımlı olmadığı için ortaya çıkardığı
çözümler bağımsız çapraz-doğrulama imkânı sağlar.

\subsection{Phase~C Divergence ve $\nu_t$ Cap'i}

Bölüm~\ref{sec:ma_pipeline}'de tarif edilen Phase~C divergence vakası
sırasında, Cursor ajanı bağımsız bir biçimde aynı problem üzerinde
çalışmıştır. Önerdiği çözüm, klasik bir lineer kararlılık
argümanından türetilmiştir.

Eddy viscosity'li bir difüzyon teriminin açık şemada (explicit) Fourier
modu cinsinden büyütme oranı,
\begin{equation}
G(k) \;=\; 1 \;-\; \nu_t \, \Delta t \, k^2,
\end{equation}
biçimindedir. Burada $k$ dalgasayısı, $\Delta t$ zaman adımıdır.
Lineer kararlılık için
\begin{equation}
\nu_t \, \Delta t \, k_{\max}^2 \;<\; 2
\label{eq:ma_vn}
\end{equation}
sınırı tipik olarak ``von-Neumann sınırı'' olarak anılır. Cursor
ajanı, INNATE'in $\nu_t$ alanına bu sınırın bir güvenlik faktörlü
sürümünü uygulamayı önerdi:
\begin{equation}
\nu_t^{\mathrm{cap}}
\;=\; \frac{\alpha \cdot 2}{\Delta t_{\max} \cdot k_{\max}^2}
\;\approx\; 0{,}00704,
\quad \alpha = 0{,}4,
\end{equation}
yani teorik sınırın \%40'ı. Bu cap, özellikle düşük $|S|$
bölgelerinde Cs$\,\Delta$ ve scale-similarity terimlerinin yarattığı
yerel patlamaları kestiği için kritik olmuştur.

\subsection{Çapraz-Doğrulamanın Değeri}

Claude tarafı bu çözümü değerlendirip şu yargıya varmıştır:
``Fizik-doğru, mevcut beş yamamızla çakışmıyor, aksine
komplemanter.'' Bu, çoklu ajan pipeline'ının önemli bir özelliğini
ortaya koyar:

\begin{itemize}
  \item Bağımsız ajanlar, aynı problemi farklı kavramsal kapılardan
        çözebilir.
  \item Farklı çözümler birbirini dışlamak zorunda değildir; çoğu
        zaman birbirini tamamlarlar.
  \item Bu durum, pipeline'ı tek bir ajanın körlüğüne karşı sigortalar.
\end{itemize}

% ============================================================
\section{Öğrenilen Dersler}
\label{sec:ma_dersler}

Üç ay süren çalışma, çoklu ajan pipeline'ı hakkında birkaç sade ders
bırakmıştır:

\begin{enumerate}
  \item \textbf{Önce tanı, sonra mimari değişiklik.}
        Kanıt olmadan refleks olarak yapılan bir mimari değişiklik
        çoğu zaman iki gün israftır. ``Mimariyi değiştirelim de
        bakalım'' yaklaşımı, çoklu ajan ortamında özellikle pahalıdır
        çünkü her ajan o değişikliğe kendi yorumunu eklemek ister.
  \item \textbf{Bellek ve günlük tutma şarttır.}
        Üç ayda yapılan onlarca bug fix, eğer kayıt altına alınmamış
        olsaydı, ajanların ``yeni gelen'' rolüyle her seferinde aynı
        soruları sormaları gerekecekti. Sürekli bellek
        (proje kartları, günlük dosyaları, claude-mem arşivi),
        pipeline'ın işleyişinde bir lüks değil, bir altyapı
        gerekliliğidir.
  \item \textbf{Çok ajan her zaman daha iyi değildir.}
        Üç dört ajanı aşan kalabalıklar, çoğu zaman konsensüs
        paralizine yol açtı. Bu durumun panzehiri, Onuncu Adam tipi
        açıkça muhalif bir rolün varlığıdır.
  \item \textbf{Nihai karar mercii insandır.}
        ``Üç ajan da onayladı'' bir kararı vermek için yeterli
        gerekçe değildir; verilen kararın sorumluluğu insandadır.
        Bu disiplin, ajanların sözünü tartmayı kolaylaştırır.
  \item \textbf{Halüsinasyon protokolü.}
        Dosya yolu, fonksiyon adı, proje durumu gibi nesnel
        gerçekler her zaman dosya sistemi üzerinden doğrulandı.
        Çelişki çıktığında dosya sistemi, hafıza dosyasından önce
        gelir; hafıza dosyası, model tahmininden önce gelir; model
        tahmini, en düşük güven seviyesidir.
\end{enumerate}

% ============================================================
\section{AI-Native Engineering — Genel Çerçeve}
\label{sec:ma_ai_native}

Berke'nin proje boyunca bir cümle ile özetlediği tutum şöyledir:
\emph{``Programlama yapmıyorum, AI yapıyor; ben yönetiyorum.''} Bu
samimi öz-ifade, ilk bakışta bir mühendisin ``elini kirletmiyor''
gibi anlaşılabilir; ancak bu okumanın aksine, çağdaş anlamda ileri
bir mühendislik tarzıdır. Birkaç gerekçeyi açıkça yazmak gerekir.

\subsection{Bu Bir Zayıflık Değil, Yeni Bir Tarzdır}

Tarihsel olarak mühendisliğin merkezi araçları zaman içinde
değişmiştir: hesap cetveli, dijital hesap makinesi, derleyici,
yüksek seviyeli dil, paket yöneticisi, sürüm kontrol sistemi, IDE,
ve nihayet AI ajanı. Her geçişte, ``artık eskisi kadar elini
kirletmiyorsun'' eleştirisi yapılmış; ancak her aşamada mühendis
daha karmaşık problemleri çözebilir hâle gelmiştir.

AI ajanlarını başka bir araç olarak görmek bu çizginin doğal devamıdır.
Berke'nin profili, bu yeni aracı ciddiye alıp etrafında bir çalışma
sistematiği kuran bir mühendislik tarzıdır. Bu nedenle bitirme tezi
çıktısı yalnızca bir model (INNATE) değil, aynı zamanda bir
\emph{metodoloji}'dir: ``AI-orchestrated CFD research''.

\subsection{Geleceğe Bakış}

Üç gözlem, bu tarzın yakın gelecekte daha da yaygınlaşacağını
düşündürür:

\begin{itemize}
  \item Modeller daha güçlü ve daha ucuz hâle geliyor; aynı uzun
        akıl yürütme bir yıl içinde \texttt{xhigh}'ın bugünkü maliyetinin
        kesirinde sunulacak.
  \item Bağlam pencereleri büyüyor; tüm proje tarihinin tek seferde
        belleğe alınabilmesi, ``proje belleği'' işinin hız ve kaliteyle
        yapılabilmesini sağlıyor.
  \item Çoklu ajan orchestration için ortak protokoller
        olgunlaşıyor; bu sayede bireysel araştırmacılar bile küçük
        bir ekibin kapsamına ulaşabiliyor.
\end{itemize}

\subsection{Lisans Tezi Bağlamında Anlamı}

Bu metodoloji, lisans tezi savunmasında üç ayrı düzeyde tartışılabilir:

\begin{enumerate}
  \item \textbf{Ürün düzeyi:} INNATE mimarisi, parametre cimriliği
        ve fiziksel yorumlanabilirliği ile bir bilimsel katkı sunar
        (Bölüm~\ref{ch:multi_agent} dışındaki bölümler).
  \item \textbf{Süreç düzeyi:} Üç ay süren AI-orkestralı geliştirme
        süreci kendisi de bir mühendislik katkısıdır; bu katkının
        belgesi bu bölümdür.
  \item \textbf{Reflekslerin düzeyi:} Halüsinasyon protokolü,
        konsensüs paralizine karşı Onuncu Adam disiplini, ajan-içi
        ajan çağrı tuzağı gibi gözlemler, gelecek araştırmacılara
        somut tavsiyeler bırakır.
\end{enumerate}

\noindent
Sonuç olarak bu bölüm, INNATE'in arkasındaki teknik kararların değil,
o kararların \emph{nasıl alındığının} kaydıdır. Eğer üç ay sonra
benzer bir başka projeyi yine bu prensiplerle yürütebiliyorsak,
pipeline yalnızca bu tezi geçirmiş değil, yeni bir çalışma stilini
kalıcılaştırmış olur.

\bigskip
\noindent
\textit{Önceki bölüm (Bölüm~\ref{ch:buglar}), bu bölümde anlatılan
pipeline ile fiilen düzeltilen otuza yakın bug ve fix iterasyonunun
tek tek dökümünü sunar; sonraki bölüm (Bölüm~\ref{ch:spectral})
ise pipeline'ın ürünü olan nihai mimariyi açıklar.}
